% Sezione 6: Sviluppi Futuri
\section{Conclusioni e Sviluppi Futuri}

\begin{frame}
    \frametitle{Risultati Raggiunti}
    \begin{block}{Progettazione Completa AMR}
        \begin{itemize}
            \item Requisiti: 1000 kg capacità, 1 m/s, 4h autonomia, €15k budget
            \item Configurazione 2 motorwheels + 2 caster (diagonali)
            \item Componenti: CFR MRT05, FlashBattery LiFePO4, Livox Mid-360, LattePanda Sigma
            \item Verifiche energetiche: motori 18\% margine piano, batteria 7.5\% margine
        \end{itemize}
    \end{block}

    \begin{block}{Modellazione e Controllo}
        \begin{itemize}
            \item Modello Bicycle: singolarità su crab/spin (FAIL)
            \item Modello Bisteer: vincoli non-olonomi, least-squares, nessuna singolarità
            \item Simulazioni: RMSE < 1.5 cm su 4 traiettorie (straight, arc, crab, spin)
            \item PID+Feedforward: riduzione RMSE ~30\%, settling time < 2s
        \end{itemize}
    \end{block}

    \begin{block}{Integrazione Tecnologie}
        \begin{itemize}
            \item FAST-LIO 2 (iEKF + ikd-Tree) per SLAM > 100 Hz
            \item CANopen CiA 402 per controllo deterministic motori (PDO, heartbeat)
        \end{itemize}
    \end{block}
\end{frame}

\begin{frame}
    \frametitle{Sviluppi Futuri}
    \begin{enumerate}
        \item \textbf{Implementazione Software}
              \begin{itemize}
                  \item Integrazione ROS 2 per architettura modulare
                  \item Deployment algoritmo FAST-LIO 2 su LattePanda Sigma
                  \item Interfaccia CANopen per controllo motori in tempo reale
              \end{itemize}

        \item \textbf{Validazione Sperimentale}
              \begin{itemize}
                  \item Test su campo presso stabilimento Daikin Ariccia
                  \item Misurazione prestazioni: precisione tracking, autonomia, affidabilità
                  \item Ottimizzazione parametri controllo su dati reali
              \end{itemize}

        \item \textbf{Estensioni}
              \begin{itemize}
                  \item Fleet management per coordinamento multi-robot
                  \item Integrazione con MES aziendale per pianificazione trasporti
                  \item Sensori addizionali: telecamere RGB-D, safety laser scanner
              \end{itemize}
    \end{enumerate}
\end{frame}
\begin{frame}
    \frametitle{Conclusioni e Sviluppi Futuri}
    \begin{columns}
        \column{0.5\textwidth}
        \begin{block}{Stato Attuale}
            \begin{itemize}
                \item Definizione requisiti e componenti completata.
                \item Analisi cinematica e scelta algoritmi (FAST-LIO 2) effettuata.
                \item Acquisizione hardware in corso.
            \end{itemize}
        \end{block}

        \column{0.5\textwidth}
        \begin{block}{Prossimi Passi}
            \begin{enumerate}
                \item \textbf{Assemblaggio}: Costruzione telaio e cablaggio elettrico.
                \item \textbf{Integrazione Software}: Setup ROS 2 e driver CANopen.
                \item \textbf{Test in Campo}:
                \begin{itemize}
                    \item Mapping dello stabilimento di Ariccia.
                    \item Test di navigazione con ostacoli dinamici.
                \end{itemize}
                \item \textbf{Integrazione WMS}: Interfacciamento con gestionale di magazzino.
            \end{enumerate}
        \end{block}
    \end{columns}
\end{frame}
